% \documentclass[12pt]{article}

%%v helpful MIT exam schema, see http://www-math.mit.edu/~psh/exam/examdoc.pdf
 \documentclass[9pt]{exam} 
% \qformat{\textbf{Question \thequestion}\quad \thepoints\hfill } % Large depth to make space}
 \renewcommand{\thesubpart}{(\roman{subpart})}
 \renewcommand{\subpartlabel}{\thesubpart}
 \footer{}{}{Page \thepage\ de \numpages}
\addpoints
% \boxedpoints
\renewcommand{\questionshook}{\setlength{\itemsep}{0.2cm}}
\renewcommand{\partshook}{
	\setlength{\itemsep}{0.1cm}
}


\usepackage{cmbright}

\pagestyle{empty} %for handouts
\addtolength{\jot}{2em} % this makes all formulae a bit more spaced vertically for handouts
\setlength{\parindent}{0pt} %we're not writing a novel
%\renewcommand{\arraystretch}{2} %makes tables taller, so easier to write in

\usepackage[fleqn]{amsmath}
\usepackage{amssymb} %standard classy maths symbols
\usepackage[a4paper, margin=1in]{geometry} %makes it easy to specify where to put stuff on the page

\usepackage{siunitx} %v handy way of getting SI units to typeset correctly...
% \sisetup{locale = FR} %... and now it will use a comma for decimals

\usepackage[utf8]{inputenc}

% \usepackage[french]{babel} % frenchify everything (e.g. a Table becomes a Tableau)
% \usepackage{lmodern} % font with nicer French characters than standard
% \usepackage{textcomp} % and more help for special characters

\usepackage{tcolorbox} % basic but good package for boxes around text
\tcbset{colback=black!5!white}
% \usepackage{color} % colours
%\newcommand{\fillin}{\color{black!1!white}} %makes the text disappear
%\definecolor{fillin}{rgb} {1.00,1.00,1.00}
%\renewcommand{\fillin}{\color{blue}} \definecolor{fillin}{rgb} {0.00,0.00,1.00} %makes the filling in text appear (in blue)

\usepackage{enumitem} % can change the labels on lists, items, enumerates
\usepackage{setspace}

\usepackage{tikz, pgfplots} % interval diagrams, sets, etc.
\usetikzlibrary{calc,trees,positioning,arrows,fit,shapes}

\begin{document}

We will study the function $\displaystyle f(x) = \frac{1}{x+2} - \frac{1}{x+3}$ throughout. \quad \textbf{Name:} \dotfill

\begin{questions}

\question 
\begin{parts}
\part[2]
Determine the domain of $f$.
\fillwithdottedlines{15mm}
\part[1] Calculate the $y$-intercept of $f$. \hfill \emph{As a reduced fraction! Without that calculator!}
\fillwithdottedlines{15mm}
\part[3] Simplify the expression to show that $\displaystyle f(x) = \frac{1}{(x+2)(x+3)}$
\fillwithdottedlines{50mm}
\part[1] Give the zeros of $f$.
\fillwithdottedlines{15mm}
\part[1] Explain how you can tell that $f$ has a horizontal asymptote $\{y = 0\}$.
\fillwithdottedlines{10mm}

\end{parts}

\begin{minipage}[t]{0.4\textwidth}
\question[5] Calculate the sign table for $f$.

\

\renewcommand{\arraystretch}{2}
\begin{tabular}{|p{15mm}|*{5}{p{5mm}|}}
\hline
&  &  &  &  &  \\ \hline
    &  &  &  &  &  \\ \hline
    &  &  &  &  &  \\ \hline
    &  &  &  &  &  \\ \hline
\end{tabular}

\end{minipage}
\hfill 
\begin{minipage}[t]{0.4\textwidth}
\question[5] Sketch the curve $y=f(x)$. 

\

\begin{tikzpicture}

\draw[-latex] (-5, 0) -- (2, 0);
\draw[-latex] (0, -2) -- (0, 2);

\end{tikzpicture}

\end{minipage}

\vfill
\question \textbf{Bonus.} Show that $x^n+1$  is divisible by $x+1$ if, and only if, $n$ is an odd integer.


\end{questions}

\begin{tcolorbox}

\textbf{Name of marker:}

\begin{center}
\gradetable[h][questions]
\end{center}

\end{tcolorbox}


\end{document}